% GENERATED FILE. Edit canonical lessons, then run npm run generate:books.
% canonical-source-hash: fnv1a64:004fda56652facb9
% canonical-lessons: MR-C117-ghadyal, MR-C117-topi, MR-C117-but, MR-C117-chappal, MR-C117-tis

\chapter{Watch, Cap, Shoe, Sandal, Thirty}
\label{ch:mr-watch-cap-shoe-sandal-thirty}
\hlchaptermodality{117}

\begin{chapteropening}
\textbf{By the end of this chapter:} \emph{I can say a watch, a clock, a cap, a shoe, a boot, a sandal and thirty.}

Complete the last lesson of chapter 117: I can say a watch, a clock, a cap, a shoe, a boot, a sandal and thirty.
\end{chapteropening}

\section[\texorpdfstring{\mr{घड्याळ} (ghaḍyāḷ)}{घड्याळ (ghaḍyāḷ)}]{\mr{घड्याळ} (ghaḍyāḷ) — a watch, a clock}

\label{lesson:MR-C117-ghadyal}



\begin{quote}
Before the new one: say the Marathi for a pocket, then the Marathi for a ring.
\end{quote}

\subsection*{You'll want to know: \mr{घड्याळ}}
\textbf{\mr{घड्याळ}} — \emph{ghaḍyāḷ} — \textquotedblleft{}a watch, a clock\textquotedblright{}.

\begin{cousinweb}
\textbf{\mr{घड्याळ}} (\emph{ghaḍyāḷ}) is a tadbhava, a word inherited and reshaped by sound change. It comes from \emph{*ghaṭītāḍa}, \textquotedblleft{}striking the time\textquotedblright{}, a form CDIAL reconstructs from its descendants (CDIAL 4413). It is \emph{ghaṭī}, \textquotedblleft{}water-clock\textquotedblright{}, and \emph{tāḍa}, \textquotedblleft{}blow\textquotedblright{}: the gong struck to mark the hours.
\end{cousinweb}

\begin{grammarlens}[title={What you've built}]
One more word for this chapter.
\end{grammarlens}

\subsection*{Your turn}
\begin{itemize}
\raggedright
  \item \emph{Say it:} \emph{ghaḍyāḷ}
  \item \emph{Say it:} \emph{ghaḍyāḷ}, once more
  \item \emph{Say it:} say \emph{ghaḍyāḷ}
  \item \emph{Recall:} say the Marathi for a pocket, then the Marathi for a ring, then say \emph{ghaḍyāḷ} again
\end{itemize}

\begin{tcolorbox}[breakable,colback=teal!4,colframe=teal!35!black,title={Before you move on}]
What does \mr{घड्याळ} mean? (\textquotedblleft{}A watch, a clock\textquotedblright{}.) Say it once more.
\end{tcolorbox}
\section[\texorpdfstring{\mr{टोपी} (ṭopī)}{टोपी (ṭopī)}]{\mr{टोपी} (ṭopī) — a cap}

\label{lesson:MR-C117-topi}



\begin{quote}
Before the new one: say the Marathi for a ring, then the Marathi for a watch.
\end{quote}

\subsection*{You'll want to know: \mr{टोपी}}
\textbf{\mr{टोपी}} — \emph{ṭopī} — \textquotedblleft{}a cap\textquotedblright{}.

\begin{cousinweb}
\textbf{\mr{टोपी}} (\emph{ṭopī}) is a tadbhava, a word inherited and reshaped by sound change. It comes from \emph{*ṭōppa}, \textquotedblleft{}hat, covering\textquotedblright{}, a form CDIAL reconstructs (CDIAL 5481). CDIAL thinks a link with Sanskrit \emph{āṭopa} very doubtful.
\end{cousinweb}

\begin{grammarlens}[title={What you've built}]
One more word for this chapter.
\end{grammarlens}

\subsection*{Your turn}
\begin{itemize}
\raggedright
  \item \emph{Say it:} \emph{ṭopī}
  \item \emph{Say it:} \emph{ṭopī}, once more
  \item \emph{Say it:} say \emph{ṭopī}
  \item \emph{Recall:} say the Marathi for a ring, then the Marathi for a watch, then say \emph{ṭopī} again
\end{itemize}

\begin{tcolorbox}[breakable,colback=teal!4,colframe=teal!35!black,title={Before you move on}]
What does \mr{टोपी} mean? (\textquotedblleft{}A cap\textquotedblright{}.) Say it once more.
\end{tcolorbox}
\section[\texorpdfstring{\mr{बूट} (būṭ)}{बूट (būṭ)}]{\mr{बूट} (būṭ) — a shoe, a boot}

\label{lesson:MR-C117-but}



\begin{quote}
Before the new one: say the Marathi for a watch, then the Marathi for a cap.
\end{quote}

\subsection*{You'll want to know: \mr{बूट}}
\textbf{\mr{बूट}} — \emph{būṭ} — \textquotedblleft{}a shoe, a boot\textquotedblright{}.

\begin{cousinweb}
\textbf{\mr{बूट}} (\emph{būṭ}) is a loan from English \emph{boot}.
\end{cousinweb}

\begin{grammarlens}[title={What you've built}]
One more word for this chapter.
\end{grammarlens}

\subsection*{Your turn}
\begin{itemize}
\raggedright
  \item \emph{Say it:} \emph{būṭ}
  \item \emph{Say it:} \emph{būṭ}, once more
  \item \emph{Say it:} say \emph{būṭ}
  \item \emph{Recall:} say the Marathi for a watch, then the Marathi for a cap, then say \emph{būṭ} again
\end{itemize}

\begin{tcolorbox}[breakable,colback=teal!4,colframe=teal!35!black,title={Before you move on}]
What does \mr{बूट} mean? (\textquotedblleft{}A shoe, a boot\textquotedblright{}.) Say it once more.
\end{tcolorbox}
\section[\texorpdfstring{\mr{चप्पल} (chappal)}{चप्पल (chappal)}]{\mr{चप्पल} (chappal) — a sandal}

\label{lesson:MR-C117-chappal}



\begin{quote}
Before the new one: say the Marathi for a cap, then the Marathi for a shoe.
\end{quote}

\subsection*{You'll want to know: \mr{चप्पल}}
\textbf{\mr{चप्पल}} — \emph{chappal} — \textquotedblleft{}a sandal\textquotedblright{}.

\begin{cousinweb}
CDIAL has no entry for \textbf{\mr{चप्पल}} (\emph{chappal}). DEDR compares Hindi \emph{cappal}, with a \textquotedblleft{}Cf.\textquotedblright{}, to Dravidian words for \textquotedblleft{}sandal\textquotedblright{}, such as Tamil \emph{ceruppu} and Telugu \emph{ceppu} (DEDR 1963).
\end{cousinweb}

\begin{grammarlens}[title={What you've built}]
One more word for this chapter.
\end{grammarlens}

\subsection*{Your turn}
\begin{itemize}
\raggedright
  \item \emph{Say it:} \emph{chappal}
  \item \emph{Say it:} \emph{chappal}, once more
  \item \emph{Say it:} say \emph{chappal}
  \item \emph{Recall:} say the Marathi for a cap, then the Marathi for a shoe, then say \emph{chappal} again
\end{itemize}

\begin{tcolorbox}[breakable,colback=teal!4,colframe=teal!35!black,title={Before you move on}]
What does \mr{चप्पल} mean? (\textquotedblleft{}A sandal\textquotedblright{}.) Say it once more.
\end{tcolorbox}
\section[\texorpdfstring{\mr{तीस} (tīs)}{तीस (tīs)}]{\mr{तीस} (tīs) — thirty}

\label{lesson:MR-C117-tis}



\begin{quote}
Before the new one: say the Marathi for a shoe, then the Marathi for a sandal.
\end{quote}

\subsection*{You'll want to know: \mr{तीस}}
\textbf{\mr{तीस}} — \emph{tīs} — \textquotedblleft{}thirty\textquotedblright{}.

\begin{cousinweb}
\textbf{\mr{तीस}} (\emph{tīs}) is a tadbhava, a word inherited from Sanskrit and reshaped by sound change. It comes from \emph{triṁśat}, \textquotedblleft{}thirty\textquotedblright{} (CDIAL 6015).
\end{cousinweb}

\begin{grammarlens}[title={What you've built}]
One more word for this chapter.
\end{grammarlens}

\subsection*{Your turn}
\begin{itemize}
\raggedright
  \item \emph{Say it:} \emph{tīs}
  \item \emph{Say it:} \emph{tīs}, once more
  \item \emph{Say it:} say \emph{tīs}
  \item \emph{Recall:} say the Marathi for a shoe, then the Marathi for a sandal, then say \emph{tīs} again
\end{itemize}

\begin{tcolorbox}[breakable,colback=teal!4,colframe=teal!35!black,title={Before you move on}]
What does \mr{तीस} mean? (\textquotedblleft{}Thirty\textquotedblright{}.) Say it once more.
\end{tcolorbox}
